\documentclass[11pt]{article}
\usepackage{manual_docs_style}
\externaldocument{build/terminology_shared_utilities}[terminology_shared_utilities.pdf]
\externaldocument{build/environment_profiles}[environment_profiles.pdf]
\externaldocument{build/run_graph_schema}[run_graph_schema.pdf]

\begin{document}

\docTitle{Generation Policy Schema}
\phantomsection\label{docstart:generation_policy_schema}

\section*{Overview}
A generation policy defines:

\begin{itemize}
    \item how many candidate families to generate;
    \item how to sample baseline parameters and initial states;
    \item how to sample intervention times;
    \item how to sample post-intervention values;
    \item how many children to generate per root-cause parameter and direction;
    \item which validation profile generated candidates should receive.
\end{itemize}

Generation policies live under:

\begin{DocCode}
models/simulink/<MODEL>/generation_policies/<policy_id>.json
\end{DocCode}

They are consumed by the candidate-generation stage, which resolves the policy into a concrete run graph:

\begin{DocCode}
models/simulink/<MODEL>/plans/<policy_id>/run_nodes.jsonl
models/simulink/<MODEL>/plans/<policy_id>/run_edges.jsonl
models/simulink/<MODEL>/plans/<policy_id>/generation_report.json
models/simulink/<MODEL>/plans/<policy_id>/validation_report.json
\end{DocCode}



\paragraph{Acceptance criteria.}
A generated plan is accepted only if it satisfies the policy's acceptance criteria.
These criteria should be checked by the candidate-generation stage and summarized in \code{validation_report.json} and \code{generation_report.json}.

Typical acceptance criteria include:
\begin{itemize}
    \item all generated run nodes pass schema validation;
    \item all run IDs and recipe hashes are unique;
    \item all sampled values lie inside \code{experiment_config.json::allowed_intervals};
    \item every intervention changes exactly one root-cause parameter;
    \item every intervention has a matched no-intervention baseline and time-zero/static-change run;
    \item the number of generated candidate families is at least \code{target_candidate_families};
    \item skipped parameter-direction requests are recorded in \code{generation_report.json};
    \item coverage targets specified by the policy are either satisfied or explicitly reported as unmet.
\end{itemize}

Note that the generation policy does not guarantee that generated intervention trajectories are detectable or distinguishable after simulation.
Those properties are evaluated by the metric and oracle stages after the runs have been materialized,

\section*{Top-Level Schema}

A generation policy is a JSON object with the following top-level fields:

\begin{itemize}
    \item \code{policy_id}: non-empty string identifying the policy.
    \item \code{model}: model name under \code{models/simulink/<MODEL>/}.
    \item \code{seed}: integer seed used to make candidate generation reproducible.
    \item \code{target_candidate_families}: It defines how many baseline families to sample.
    \item \code{baseline_sampler}: object defining how baseline parameters and initial states are sampled.
    \item \code{intervention_sampler}: object defining intervention times, parameters, and post-intervention values.
    \item \code{adaptive_generation}: optional object defining adaptive backfilling targets.
    \item \code{artifact_policy}: optional object defining whether existing generated plans may be overwritten.
\end{itemize}

Minimal shape:

\begin{DocCode}
{
  "policy_id": "benchmark_v1",
  "model": "BallDrop",
  "seed": 123,
  "target_candidate_families": 1000,
   "baseline_sampler": {
    "method": "latin_hypercube"
    "..."
  },
  "intervention_sampler": {
    "children_per_parameter_direction": 3,
    "intervention_time": {
      "distribution": "uniform",
      "min_fraction_of_horizon": 0.2,
      "max_fraction_of_horizon": 0.8
    },
  }
}
\end{DocCode}



\section*{\code{baseline_sampler}}

\code{baseline_sampler} defines how baseline parameters and initial-state variables are sampled.
It samples over the variables listed in:

\begin{DocCode}
experiment_config.json::exposed_variables.parameters
experiment_config.json::exposed_variables.initial_state
\end{DocCode}

The sampler must respect all \code{allowed_intervals} defined in \code{experiment_config.json}.

Schema:

\begin{DocCode}
"baseline_sampler": {
  "method": "latin_hypercube",
  "variables": {
    "mass": {
      "distribution": "loguniform"
    },
    "gravity": {
      "distribution": "uniform"
    },
    "initial_height": {
      "distribution": "uniform"
    }
  },
  "baseline_min_distance_scale": 0.2
}
\end{DocCode}

Fields:

\begin{itemize}
    \item \code{method}: required string. How the joint baseline samples are generated.
    \item \code{variables}: optional object keyed by exposed variable name. It overrides the default distribution for selected variables.
    \item \code{baseline_min_distance_scale}: optional non-negative number. If present, generated baselines should be separated from one another by at least this fraction of the environment-level \termMinAbsDist{} and \termMinSRDDistance{} values for at least one exposed variable.
\end{itemize}

Allowed \code{method} values:

\begin{itemize}
    \item \code{independent}: sample each variable independently.
    \item \code{latin_hypercube}: use Latin-hypercube sampling over exposed variables.
    \item \code{sobol}: use a Sobol sequence over exposed variables.
    \item \code{grid}: generate a small Cartesian or quasi-Cartesian diagnostic grid.
\end{itemize}

Allowed variable-level \code{distribution} values:

\begin{itemize}
    \item \code{uniform}: sample uniformly in the allowed interval.
    \item \code{loguniform}: sample log-uniformly in the allowed interval. The allowed interval must be strictly positive.
    \item \code{boundary_biased}: sample with extra mass near the lower and upper legal bounds.
    \item \code{fixed}: use a fixed value. Requires a \code{value} field.
\end{itemize}

Example fixed variable:

\begin{DocCode}
"variables": {
  "gravity": {
    "distribution": "fixed",
    "value": 9.81
  }
}
\end{DocCode}

A fixed value must still lie inside the variable's \code{allowed_intervals}.

\section*{\code{intervention_sampler}}

\code{intervention_sampler} defines how intervention children are generated from each baseline family.

Schema:

\begin{DocCode}
"intervention_sampler": {
  "parameters": "all",
  "directions": ["increase", "decrease"],
  "children_per_parameter_direction": 3,
  "intervention_time": {
    "distribution": "uniform",
    "min_fraction_of_horizon": 0.2,
    "max_fraction_of_horizon": 0.8
  },

}
\end{DocCode}

Fields:

\begin{itemize}
    \item \code{parameters}: either \code{"all"} or a non-empty list of parameter names from \code{experiment_config.json::exposed_variables.parameters}.
    \item \code{directions}: non-empty list containing \code{"increase"}, \code{"decrease"}, or both.
    \item \code{children_per_parameter_direction}: positive integer.
    \item \code{intervention_time}: object defining how intervention times are sampled.
\end{itemize}

Only variables listed in \code{experiment_config.json::exposed_variables.parameters} may be intervened on.
Variables listed under \code{initial_state} are baseline-only variables and must not be changed after simulation start.

\subsection*{\code{intervention_time}}

\code{intervention_time} defines how the step-change time is sampled.

Recommended schema:

\begin{DocCode}
"intervention_time": {
  "distribution": "uniform",
  "min_fraction_of_horizon": 0.2,
  "max_fraction_of_horizon": 0.8
}
\end{DocCode}

Allowed fields:

\begin{itemize}
    \item \code{distribution}: required string. Allowed values are \code{uniform}, \code{fixed}, and \code{choice}.
    \item \code{min_fraction_of_horizon}: optional float in \code{[0,1]}.
    \item \code{max_fraction_of_horizon}: optional float in \code{[0,1]}.
    \item \code{value}: required when \code{distribution} is \code{fixed}.
    \item \code{values}: required when \code{distribution} is \code{choice}.
\end{itemize}

All intervention times must satisfy:

\begin{DocCode}
0 < intervention_time < experiment_config.json::end_time_input_s
\end{DocCode}

If \code{min_fraction_of_horizon} and \code{max_fraction_of_horizon} are used, the actual time interval is computed from \code{experiment_config.json::end_time_input_s}.

\subsection*{\code{value_sampler}}

\code{value_sampler} defines how the post-intervention value is sampled for each changed parameter.

Recommended schema:

\begin{DocCode}
"value_sampler": {
  "method": "conditional_uniform",
  "min_abs_distance_multiplier": 1.0,
  "min_srd_distance_multiplier": 1.0
}
\end{DocCode}

Allowed \code{method} values:

\begin{itemize}
    \item \code{conditional_uniform}: uniformly sample valid values on the requested side of the baseline value after applying distance constraints.
    \item \code{boundary_biased}: sample valid values with extra mass near the legal interval boundaries.
    \item \code{fixed_delta}: add or subtract a fixed amount. Requires \code{delta}.
    \item \code{fixed_srd}: sample values targeting a fixed SRD distance. Requires \code{target_srd}.
\end{itemize}

The generated post-intervention value must satisfy all of the following:

\begin{itemize}
    \item it lies inside \code{experiment_config.json::exposed_variables.parameters::<parameter>::allowed_intervals};
    \item it changes exactly one root-cause parameter;
    \item its absolute distance from the baseline value is at least the environment-level \termMinAbsDist{} multiplied by \code{min_abs_distance_multiplier};
    \item its \termSRD{} distance from the baseline value is at least the environment-level \termMinSRDDistance{} multiplied by \code{min_srd_distance_multiplier};
    \item its direction agrees with the requested direction.
\end{itemize}

If no valid value exists for a requested parameter and direction, the generator must not create a runnable child node.
The missing direction should instead be recorded in \code{generation_report.json}.

\section*{\code{adaptive_generation}}

\code{adaptive_generation} is optional.
It specifies whether the generator should run in batches and backfill underrepresented classes or directions after eligibility computation.

Schema:

\begin{DocCode}
"adaptive_generation": {
  "enabled": true,
  "target_eligible_children_per_parameter": 200,
  "target_eligible_baselines": 100,
  "batch_size_families": 100,
  "max_batches": 20,
  "rebalance_by": ["parameter", "direction"]
}
\end{DocCode}

Fields:

\begin{itemize}
    \item \code{enabled}: boolean.
    \item \code{target_eligible_children_per_parameter}: optional positive integer.
    \item \code{target_eligible_baselines}: optional positive integer.
    \item \code{batch_size_families}: positive integer.
    \item \code{max_batches}: positive integer.
    \item \code{rebalance_by}: optional list of grouping keys. Allowed values include \code{parameter}, \code{direction}, and \code{validation_profile}.
\end{itemize}

Adaptive generation should not change the meaning of previously generated run IDs.
It should only add new candidate families or candidate children.

\section*{\code{manual_cases}}

\code{manual_cases} is optional.
It specifies whether curated human-authored cases should be merged into the resolved run graph.

Schema:

\begin{DocCode}
"manual_cases": {
  "enabled": true,
  "path": "manual_cases.jsonl",
  "default_validation_profile": "diagnostic"
}
\end{DocCode}

Fields:

\begin{itemize}
    \item \code{enabled}: boolean.
    \item \code{path}: path relative to \code{models/simulink/<MODEL>/}.
    \item \code{default_validation_profile}: validation profile assigned to manual cases that do not specify one explicitly.
\end{itemize}

Manual cases and programmatically generated cases are resolved into the same run graph.
They share the same run-node schema.
The distinction is recorded in node metadata:

\begin{DocCode}
"metadata": {
  "source": "manual",
  "validation_profile": "diagnostic"
}
\end{DocCode}

A human-authored case is not automatically relaxed.
Core semantic validity rules still apply.
For example, values must be finite, legal, and runnable.

\section*{\code{artifact_policy}}

\code{artifact_policy} is optional.
It defines how generated artifacts are written.

Schema:

\begin{DocCode}
"artifact_policy": {
  "overwrite_existing_plan": false,
  "allow_append": true,
  "write_compat_model_run_specs": false
}
\end{DocCode}

Fields:

\begin{itemize}
    \item \code{overwrite_existing_plan}: whether the generator may overwrite an existing \code{plans/<policy_id>/} directory.
    \item \code{allow_append}: whether the generator may add new nodes to an existing plan without deleting old nodes.
    \item \code{write_compat_model_run_specs}: whether to emit a legacy \code{model_run_specs.json} compatibility artifact.
\end{itemize}

New code should consume the resolved run graph directly.
Legacy compatibility artifacts should be used only during migration.

\section*{Resolved Output}

A generation policy is not consumed directly by Stage: Simulate.
It is first resolved into a concrete run graph.

The resolver writes:

\begin{itemize}
    \item \code{run_nodes.jsonl}: flat list of runnable run recipes.
    \item \code{run_edges.jsonl}: relationships among baseline, intervention, and time-zero/static-change runs.
    \item \code{generation_report.json}: summary of sampled candidates, skipped directions, and coverage.
    \item \code{validation_report.json}: summary of schema and semantic validation checks.
\end{itemize}

\subsection*{\code{run_nodes.jsonl}}

Each line is one runnable run node.
The run ID is content-addressed from the canonical recipe.
Metadata does not affect the run ID.

Example:

\begin{DocCode}
{
  "run_id": "run_7fd9daa0055fb3f412ef5e0ab8d4c5ba",
  "kind": "intervention",
  "family_id": "fam_de4d7c0a5335e9005510cc2442d0971e",
  "recipe_hash": "7fd9daa0055fb3f412ef5e0ab8d4c5ba",
  "recipe": {
    "model": "BallDrop",
    "baseline_parameters": {
      "mass": 4.0,
      "drag_coeff": 0.05,
      "restitution": 0.8,
      "gravity": 9.81,
      "initial_height": 3.0,
      "initial_velocity": -2.0
    },
    "intervention": {
      "mode": "at_intervention_time",
      "parameter": "gravity",
      "value": 13.2,
      "time": 4.0
    }
  },
  "metadata": {
    "source": "programmatic",
    "policy_id": "benchmark_v1",
    "validation_profile": "benchmark_candidate",
    "seed": 123,
    "generator": "compile_run_graph.py",
    "generator_version": "1.0.0"
  }
}
\end{DocCode}

Allowed \code{kind} values:

\begin{itemize}
    \item \code{baseline}
    \item \code{intervention}
    \item \code{time0_baseline}
\end{itemize}

Allowed \code{intervention.mode} values:

\begin{itemize}
    \item \code{none}: no intervention; used for baseline runs.
    \item \code{at_intervention_time}: parameter changes at \code{intervention.time}.
    \item \code{from_time_zero}: same post-intervention value is active from time zero.
\end{itemize}

\subsection*{\code{run_edges.jsonl}}

Each line records a relationship among run nodes.

Example:

\begin{DocCode}
{
  "edge_type": "baseline_to_intervention",
  "family_id": "fam_de4d7c0a5335e9005510cc2442d0971e",
  "source_run_id": "run_baseline_...",
  "target_run_id": "run_intervention_..."
}
\end{DocCode}

Expected edge types:

\begin{itemize}
    \item \code{baseline_to_intervention}
    \item \code{intervention_to_time0_baseline}
    \item \code{same_family}
\end{itemize}

Edges are used by simulation, metrics, oracle optimization, web visualization, and question selection to recover the matched baseline and time-zero/static-change probe for each intervention.

\subsection*{\code{generation_report.json}}

\code{generation_report.json} records the outcome of candidate generation.

It should include:

\begin{itemize}
    \item timestamp;
    \item \code{policy_id};
    \item model;
    \item generator version;
    \item number of requested families;
    \item number of generated families;
    \item number of baseline nodes;
    \item number of intervention nodes;
    \item number of time-zero/static-change nodes;
    \item skipped parameter-direction requests;
    \item per-parameter and per-direction coverage statistics.
\end{itemize}

Skipped directions are recorded here, not as null child nodes in \code{run_nodes.jsonl}.

Example skipped direction:

\begin{DocCode}
{
  "family_id": "fam_de4d7c0a5335e9005510cc2442d0971e",
  "parameter": "restitution",
  "direction": "increase",
  "reason": "baseline value too close to upper allowed bound"
}
\end{DocCode}

\subsection*{\code{validation_report.json}}

\code{validation_report.json} records the validation result for the resolved plan.

It should include:

\begin{itemize}
    \item whether all run IDs are unique;
    \item whether all recipe hashes are unique;
    \item whether all recipes are finite and valid JSON;
    \item whether every recipe respects \code{allowed_intervals};
    \item whether every intervention changes exactly one root-cause parameter;
    \item whether every intervention time is inside the simulation horizon;
    \item whether every intervention has a matched baseline and time-zero/static-change run;
    \item whether all graph edges point to existing nodes.
\end{itemize}

\section*{Validation Rules}

The following rules are always enforced, regardless of \code{source} or \code{validation_profile}:

\begin{itemize}
    \item all run IDs must be unique;
    \item all recipe hashes must be unique;
    \item all numeric values must be finite;
    \item every recipe must specify a valid model;
    \item every baseline recipe must contain all exposed parameters and initial-state variables;
    \item every intervention recipe must change exactly one variable from \code{experiment_config.json::exposed_variables.parameters};
    \item no intervention may change a variable listed only under \code{initial_state};
    \item every intervention value must lie inside the parameter's \code{allowed_intervals};
    \item every intervention time must satisfy \code{0 < intervention_time < end_time_input_s};
    \item every \code{time0_baseline} must correspond to exactly one intervention node;
    \item every edge must point to an existing node.
\end{itemize}

The following rules apply to \code{benchmark_candidate} and \code{paper_selected} profiles:

\begin{itemize}
    \item intervention values must satisfy the configured \termMinAbsDist{} constraint;
    \item intervention values must satisfy the configured \termMinSRDDistance{} constraint;
    \item baseline families should provide coverage across root-cause parameters;
    \item generated baselines should be sufficiently diverse according to the policy's diversity settings;
    \item generated candidates should include matched no-intervention and time-zero/static-change probes.
\end{itemize}

The following rules are coverage goals rather than hard semantic validity rules:

\begin{itemize}
    \item every family should contain increase and decrease examples for every root-cause parameter when feasible;
    \item every parameter should have at least the target number of eligible children after metric filtering;
    \item final selected datasets should be balanced across classes.
\end{itemize}

Coverage failures should be recorded in \code{generation_report.json} or \code{selection_report.json}, not represented as fake runnable runs.

\section*{Example Policy}

\begin{DocCode}
{
  "policy_id": "benchmark_v1",
  "model": "BallDrop",
  "seed": 123,
  "validation_profile": "benchmark_candidate",

  "target_candidate_families": 1000,

  "baseline_sampler": {
    "method": "latin_hypercube",
    "variables": {
      "mass": {
        "distribution": "loguniform"
      },
      "drag_coeff": {
        "distribution": "loguniform"
      },
      "restitution": {
        "distribution": "uniform"
      },
      "gravity": {
        "distribution": "uniform"
      },
      "initial_height": {
        "distribution": "uniform"
      },
      "initial_velocity": {
        "distribution": "uniform"
      }
    },
    "baseline_min_distance_scale": 0.2
  },

  "intervention_sampler": {
    "parameters": "all",
    "directions": ["increase", "decrease"],
    "children_per_parameter_direction": 3,
    "intervention_time": {
      "distribution": "uniform",
      "min_fraction_of_horizon": 0.2,
      "max_fraction_of_horizon": 0.8
    },
    "value_sampler": {
      "method": "conditional_uniform",
      "min_abs_distance_multiplier": 1.0,
      "min_srd_distance_multiplier": 1.0
    }
  },

  "adaptive_generation": {
    "enabled": true,
    "target_eligible_children_per_parameter": 200,
    "target_eligible_baselines": 100,
    "batch_size_families": 100,
    "max_batches": 20,
    "rebalance_by": ["parameter", "direction"]
  },

  "manual_cases": {
    "enabled": true,
    "path": "manual_cases.jsonl",
    "default_validation_profile": "diagnostic"
  },

  "artifact_policy": {
    "overwrite_existing_plan": false,
    "allow_append": true,
    "write_compat_model_run_specs": false
  }
}
\end{DocCode}

\section*{Artifact Policy}

Generation policies are human-authored or agent-authored source files.
They should be committed to Git.

Compiled plans are generated artifacts.
Depending on size, \code{run_nodes.jsonl}, \code{run_edges.jsonl}, \code{generation_report.json}, and \code{validation_report.json} may be committed for reproducibility or regenerated from the policy.

The policy itself must not be edited by the generator.
If a generated plan needs to be changed, either:

\begin{itemize}
    \item edit the generation policy and regenerate the resolved plan; or
    \item add curated cases to \code{manual_cases.jsonl}; or
    \item create a new policy ID.
\end{itemize}

\section*{Recommended Command}

A typical candidate-generation command is:

\begin{DocCode}
python workflows/generate_runs/compile_run_graph.py \
  --model BallDrop \
  --policy benchmark_v1 \
  --overwrite
\end{DocCode}

The next stages should then consume the resolved plan:

\begin{DocCode}
python workflows/simulate/run_pending_sims.py \
  --model BallDrop \
  --policy benchmark_v1 \
  --compute-features

env/bin/python workflows/metrics/compute_metrics.py \
  --model BallDrop \
  --policy benchmark_v1 \
  --jobs 8
\end{DocCode}

\end{document}
